\section{Substitution coherence and open use}
\label{sec:naturality}

\subsection{Instantiation, not recompilation}

Let \(A\in\mathsf{BCode}(k)\) and let the compiler produce
\(\chi_A\in\PR\,k\).  Given arbitrary terms
\(\sigma:\mathrm{Fin}(k)\to\PTm\,n\), define
\[
  \chi_A[\sigma] := \mathsf{app}(\chi_A,\sigma),
  \qquad
  \llbracket A\rrbracket[\sigma]
     := \llbracket A\rrbracket.\mathsf{subst}(\sigma).
\]
Because the symbol is fixed, the syntax law for term application gives
\[
 \mathsf{subst}(\tau,\mathsf{app}(\chi_A,\sigma))
   = \mathsf{app}(\chi_A, i\mapsto\mathsf{subst}(\tau,\sigma_i)).
\]
There is no second compiler output to compare.

\begin{theorem}[Open reflection]
\label{thm:open-reflection}
For every context \(\Gamma\), every term substitution \(\sigma\), and every
bounded scheme \(A\), iPRA derives under \(\Gamma\)
\begin{align*}
 &\chi_A[\sigma]=0\vee\chi_A[\sigma]=1,\\
 &\bigl(\chi_A[\sigma]=1\bigr)
       \leftrightarrow\llbracket A\rrbracket[\sigma],\\
 &\bigl(\chi_A[\sigma]=0\bigr)
       \leftrightarrow\neg\llbracket A\rrbracket[\sigma].
\end{align*}
\end{theorem}

\begin{proof}
Apply the checked arbitrary-term instantiation theorem \code{IPRA.inst} to the
three closed generic derivations of \cref{thm:reflection}.  The formula
substitution and term-application simplification laws put the conclusions in
the displayed form.  Finally weaken from the empty context to \(\Gamma\).
\end{proof}

\subsection{Binders are handled once}

The denotation of a bounded body uses the standard lifted substitution under a
binder.  This law is part of the checked syntax implementation:
\[
  (\forall i\,B)[\sigma]
    = \forall i\,B[\mathsf{up}(\sigma)].
\]
The reflection compiler does not need its own binder-substitution recursion.
The body scheme has arity \(k+1\), and its characteristic symbol already takes
\((i,\vec x)\) as formal arguments.  When the outer parameters are specialized,
its argument vector becomes \((i,\sigma(\vec x))\) by ordinary application.

This is the principal technical advantage of schema-first compilation.  The
open obligation called ``arity-indexed PR abstraction and binder substitution''
in the earlier outline is split into components already present in the
repository:

\begin{itemize}
\item \code{PTm.toPR} closes each fixed term scheme to a symbol;
\item \code{recDef} permits arbitrary parameter arity;
\item \code{PTm.subst\_app} handles symbol application;
\item \code{Fm.subst\_subst}, \code{up}, and the binder lemmas handle formulas;
\item \code{IPRA.inst} transports closed derivations to arbitrary terms and
  contexts.
\end{itemize}

Only the arity-generic search proof and the structural reflection recursion
remain new object derivations.

\subsection{Naturality of a compiled proof plan}

Suppose a finite plan compiles to \(\Gamma\turnstilep G\).  For every
simultaneous substitution \(\sigma\), applying \code{Prf.subst} yields
\[
  \Gamma[\sigma]\turnstilep G[\sigma].
\]
This theorem applies to the final plan, not only to individual reflected
checks.  Therefore a theorem may be developed once at formal query parameters
and reused at trace states, least-search outputs, tuple projections, or
application-specific terms without rebuilding its proof graph.

\subsection{Universal closure}

An open theorem over parameters can be universally closed with repeated
universal introduction when the context permits it.  The framework thus does
not prove merely one theorem for each standard numeral.  It constructs a
single open object derivation, from which the ordinary universally quantified
mathematical statement follows inside iPRA.
